\chapter{Introducci\'on}

\section{Contexto y motivaci\'on}

En la era digital, el acceso a la informaci\'on nunca ha sido tan amplio ni tan fragmentado. Para una persona en Bolivia, informarse sobre la actualidad nacional implica revisar m\'ultiples medios de comunicaci\'on, abrir decenas de pesta\~nas del navegador, comparar titulares, buscar contexto en redes sociales y descartar noticias repetidas. Este proceso, que parece cotidiano e inofensivo, genera un volumen significativo de desperdicio digital.

El \textbf{desperdicio digital informativo} se manifiesta de varias formas: tiempo perdido revisando la misma noticia en distintos sitios, datos m\'oviles consumidos al cargar p\'aginas completas con publicidad y \textit{trackers}, procesamiento redundante cuando los sistemas de IA resumen art\'iculos duplicados, y exposici\'on a informaci\'on fragmentada o sin fuente verificable en redes sociales.

Esta problem\'atica se alinea directamente con los principios de \textbf{Green Tech}: promover un uso m\'as eficiente, responsable y sostenible de la tecnolog\'ia, la inteligencia artificial, el \textit{cloud}, los datos y los procesos internos. No se trata solo de ahorrar tiempo, sino de repensar c\'omo consumimos y procesamos informaci\'on digital.

\section{El problema del desperdicio digital informativo}

El desperdicio ocurre en cinco niveles interconectados:

\begin{table}[H]
\centering
\begin{tabular}{lp{10cm}}
\toprule
\textbf{Nivel} & \textbf{Problema} \\
\midrule
Usuario & Tiempo perdido revisando noticias repetidas o poco relevantes \\
Datos & Carga innecesaria de p\'aginas, im\'agenes, scripts y publicidad \\
IA & Procesamiento redundante si se resumen art\'iculos duplicados \\
Ecosistema digital & Mayor navegaci\'on, m\'as \textit{requests} y m\'as consumo de recursos \\
Confianza & Exposici\'on a contenido viral, fragmentado o sin fuente trazable \\
\bottomrule
\end{tabular}
\caption{Niveles del desperdicio digital informativo}
\end{table}

Una misma noticia puede aparecer en varios medios con t\'itulos distintos o URL diferentes. Tambi\'en puede circular como video corto, captura o publicaci\'on sin contexto en redes sociales como TikTok o Facebook. Sin una capa de deduplicaci\'on, priorizaci\'on y trazabilidad, el usuario y el sistema terminan procesando informaci\'on redundante o dif\'icil de verificar.

\section{Objetivos}

\subsection{Objetivo general}

Desarrollar e implementar una plataforma de inteligencia artificial responsable que reduzca el desperdicio digital informativo mediante la recolecci\'on, deduplicaci\'on, clasificaci\'on, priorizaci\'on y resumen eficiente de noticias locales bolivianas.

\subsection{Objetivos espec\'ificos}

\begin{enumerate}[nosep]
    \item Recolectar autom\'aticamente noticias de las principales fuentes informativas de Bolivia mediante \textit{web scraping}.
    \item Implementar un sistema de deduplicaci\'on multinivel (URL, t\'itulo y huella de contenido) para eliminar redundancias.
    \item Clasificar y priorizar noticias por categor\'ias y relevancia para el contexto nacional.
    \item Generar res\'umenes con IA aplicando criterios de eficiencia: filtrar, deduplicar y rankear antes de enviar al modelo.
    \item Medir y visualizar el impacto de la reducci\'on del flujo informativo mediante m\'etricas cuantificables.
    \item Distribuir la informaci\'on de forma personalizada seg\'un preferencias del usuario, reduciendo el consumo innecesario de datos.
\end{enumerate}

\section{Alcance y limitaciones}

\subsection{Alcance}

El proyecto abarca el desarrollo completo de una plataforma funcional que incluye:
\begin{itemize}[nosep]
    \item \textit{Web scraping} de medios de comunicaci\'on bolivianos (Radio Fides, Unitel, Red Uno, Red Bolivisi\'on, Los Tiempos, El Deber, P\'agina Siete y ATB).
    \item Pipeline de procesamiento con deduplicaci\'on, clasificaci\'on, ranking y resumen.
    \item Integraci\'on con proveedores de IA (Groq, OpenAI, GitHub Models) mediante un cliente abstracto.
    \item Aplicaci\'on web con frontend React y backend FastAPI.
    \item Base de datos PostgreSQL para persistencia.
    \item Sistema de suscripci\'on con preferencias por categor\'ia y frecuencia.
    \item Panel de m\'etricas de impacto Green Tech.
    \item Despliegue en VPS con Hostinger.
\end{itemize}

\subsection{Limitaciones}

El prototipo tiene limitaciones claramente identificadas:
\begin{itemize}[nosep]
    \item Las m\'etricas de tiempo y MB son estimaciones, no mediciones energ\'eticas directas.
    \item Los \textit{scrapers} dependen de cambios en el HTML de los medios.
    \item WhatsApp en producci\'on requiere credenciales y costos externos (Twilio).
    \item No todos los art\'iculos hist\'oricos tienen \textit{backfill} de huellas.
    \item La verificaci\'on autom\'atica de informaci\'on falsa no forma parte del MVP.
\end{itemize}
